% =============================================================================
% 2. Background -- target ~3 pp
% Four moves:
%   2.1 LLMs and recommender systems       -- classical recsys (brief), LLM recsys
%   2.2 Preference Optimization            -- DPO, IPO, SimPO, KTO with equations
%   2.3 In-Context Learning                -- ICL fundamentals, chat vs flat
%   2.4 What others have done so far       -- related work + the gap this fills
% =============================================================================

\section{Background}
\label{sec:bg}

\subsection{LLMs and Recommender Systems}
\label{sec:bg:llmrec}

Trying to guess what people like based on their previous preferences is not a new thing; recommender systems are actually one of the most researched and studied areas in computing, having created some of the most influential and widely used algorithms today \citep{koren2009matrixfact}, that power marketplaces such as Amazon, Alibaba ...

The creation of LLMs has shone a brand new light on what we are able to achieve in this area, these models have a pretty wide knowledge of the world they are able to create connections between items that would be impossible to discover in traditional systems that most times only used data from similar users or a very limited set of characteristics from the product. Furthermore, these systems have already been pre-trained so we can create these connections from very few previous user preferences. 

Even though the problem is the same, it requires a completely different approach, recommender systems mainly work in a collaborative manner, assuming the taste of a user is completely based on similarities to its peers. With LLMs it is more about steering the LLM's vast knowledge based on the prior preferences from the users. 

\subsection{Direct Preference Optimization}

Direct preference optimization is the mechanism used to align models given a pair of outputs where we know which one the user preferred, building on the earlier reinforcement-learning-from-human-feedback line of work \citep{christiano2017rlhf, ouyang2022instructgpt} but doing away with the explicit reward model. There are several types of DPOs; each one tackles the alignment problem in its own way, either for efficiency or a specific use case. The interesting differences sit in one of two design choices, the first being how you measure what's "more likely", and what you anchor the model to so it doesn't drift off and forget everything it picked up during pre-training. We compare four members of this family in this thesis: DPO, IPO, SimPO, and KTO. Three of them keep a frozen copy of the original model (the reference $\pi_{\text{ref}}$) and steer the trainable model (the policy $\pi_\theta$) relative to it; one of them drops the reference entirely. For a broader survey of the DPO landscape beyond the four methods compared here, see \citet{xiao2024dposurvey}.

\paragraph{DPO -- Direct Preference Optimization.} The Bradley-Terry-grounded baseline. For a preference pair $(x, y_w, y_l)$ with $y_w$ preferred to $y_l$:
\begin{equation}
\mathcal{L}_{\text{DPO}} = -\log \sigma\!\left(\beta\Big[\big(\log \pi_\theta(y_w \mid x) - \log \pi_{\text{ref}}(y_w \mid x)\big) - \big(\log \pi_\theta(y_l \mid x) - \log \pi_{\text{ref}}(y_l \mid x)\big)\Big]\right).
\end{equation}
The intuition is straightforward: tilt the policy toward $y_w$ relative to its prior probability, and away from $y_l$ relative to its prior probability. The reference model acts as a leash and the policy is free to move, but only so far from where it started~\citep{rafailov2023dpo}.

\paragraph{IPO -- Identity Preference Optimization.} Replaces the log-sigmoid loss with a squared loss that has a natural target margin of $1/(2\beta)$, designed to reduce overfitting under noisy preferences~\citep{azar2024ipo}:
\begin{equation}
\mathcal{L}_{\text{IPO}} = \left(h_\theta(x, y_w, y_l) - \tfrac{1}{2\beta}\right)^2, \quad h_\theta = \big(\log \pi_\theta - \log \pi_{\text{ref}}\big)\Big|_{y_w} - \big(\log \pi_\theta - \log \pi_{\text{ref}}\big)\Big|_{y_l}.
\end{equation}
The motivation is that DPO can be too aggressive: if a preference label is wrong or noisy, the log-sigmoid keeps pushing the gap wider forever, while the squared loss has a target margin it actually wants to hit and then stops. In the cold-start regime, where every label carries a lot of weight, this matters.

\paragraph{SimPO -- Simple Preference Optimization.} Reference-free; a length-normalised log-probability gap with a fixed margin $\gamma$~\citep{meng2024simpo}:
\begin{equation}
\mathcal{L}_{\text{SimPO}} = -\log \sigma\!\left(\tfrac{\beta}{|y_w|}\log \pi_\theta(y_w \mid x) - \tfrac{\beta}{|y_l|}\log \pi_\theta(y_l \mid x) - \gamma\right).
\end{equation}
SimPO throws out the reference model entirely, which roughly halves the memory and compute needed to train. It also normalises by sequence length, which matters because longer responses naturally accumulate lower log-probabilities and would otherwise dominate the gradient just because they are longer.

\paragraph{KTO -- Kahneman-Tversky Optimization.} Operates on unpaired thumbs-up / thumbs-down signals; inspired by prospect theory's asymmetric treatment of gains and losses~\citep{ethayarajh2024kto}. Per-example reward $r_\theta(x, y) = \log \pi_\theta(y \mid x) - \log \pi_{\text{ref}}(y \mid x)$; loss alternates by label:
\begin{equation}
\mathcal{L}_{\text{KTO}} =
\begin{cases}
-\log \sigma\!\big(\beta(r_\theta - z_0)\big) & \text{if } y \text{ desirable},\\
-\log \sigma\!\big(\beta(z_0 - r_\theta)\big) & \text{if } y \text{ undesirable},
\end{cases}
\end{equation}
where $z_0$ is a batch-level reference baseline. The big departure here is that KTO doesn't need preference pairs at all --- a single ``I liked this'' or ``I didn't'' is enough --- which makes it the closest match to the kind of feedback most products actually collect from users in the wild.

I chose these models because together they span the design space that matters for cold-start personalisation. DPO is the baseline that every newer method positions itself against. IPO changes the shape of the loss to be more robust to noisy or scarce data, which is exactly our setting. SimPO drops the reference, asking whether that regularisation is still worth its cost when there is very little training data to drift on. KTO drops the pairing requirement, asking whether cheaper unpaired feedback is good enough. The empirical question for this thesis is which of these trade-offs actually pays off.


\subsection{In-Context Learning}

In-context learning is the simplest approach to the alignment problem; it steers the LLM using a small set of demonstration examples placed directly in the prompt, without any fine-tuning or model update. This started being possible as models got bigger, smaller models can't really do it, but once they cross a certain size threshold they start picking up patterns from a few examples in the prompt and applying them to a new query~\citep{brown2020gpt3}. The mechanism behind this behavior was later characterized by \citet{olsson2022induction} as the formation of \emph{induction heads} during training, attention circuits that copy and complete patterns from earlier in the context. In the ICL context there are three different approaches. The first is zero-shot, where the model is given only a task description, then we have few-shot, where it is given a handful of input-output examples, and finally instruction-tuned, where the model has been trained to follow natural-language instructions and so behaves more reliably under both~\citep{dong2024iclsurvey}. For our purposes, ICL is attractive precisely because it sidesteps training altogether; the user's preferences simply become more text in the prompt.

We use two variants in this thesis. \texttt{icl\_chat} inserts the training pairs as prior context in Llama-3-Instruct's chat template, so the model sees the user's preferences as if they were earlier context of an ongoing conversation. \texttt{icl\_flat} adds the same training pairs as plain text before the main prompt, with no chat-template structure. ICL is often dismissed as a baseline because it doesn't ``learn'' by adjusting the model. This adds another layer to the analysis, instead of simply looking for the best learning method. But asking if not learning sometimes might be a better strategy?
                            
\subsection{What Others Have Done So Far}
\label{sec:bg:related}
                            
Three lines of recent work touch the question this thesis asks; none combine into the head-to-head comparison it carries out.

The first is the LLM-as-recommender literature. \citet{bao2023tallrec} introduced TALLRec, a parameter-efficient tuning framework that aligns instruction-tuned LLMs with collaborative-filtering signal and showed it could match or beat classical baselines with a fraction of the data. \citet{friedman2023recllm} built a full conversational recommender that stacks retrieval, user-profile summarisation and dialogue management on top of a frozen LLM. The recent survey by \citet{fan2024recllmsurvey} catalogues the broader space and is explicit that prompting and parameter-efficient fine-tuning are the two dominant interfaces, but it does not pit them against each other under matched data budgets.

The second is the personalised-alignment literature, which extends generic RLHF and DPO to settings where preferences differ across users. \citet{li2024prlhf} adds a user-conditional component to the reward model so a single policy can serve heterogeneous preferences. \citet{salemi2024lamp} introduces the LaMP benchmark suite for personalised generation and reports both retrieval-augmented prompting and supervised fine-tuning baselines on it. The closest existing work to this thesis is \citet{singh2025fspo}, whose FSPO method studies few-shot preference optimisation with synthetic preferences and reports transfer to real users. FSPO motivates the cold-start regime but does not compare its trained policies against in-context learning on the same preferences.

The third is the ICL literature itself, surveyed by \citet{dong2024iclsurvey} and grounded by the few-shot scaling demonstrated in \citet{brown2020gpt3}, which has documented strong cold-start behaviour across many tasks but has not been directly benchmarked against per-user preference optimisation. \citet{liu2023lostmiddle} flags one of the operational limits of the prompt route, namely that long contexts are attended to unevenly, which is part of why our experiments separate \texttt{icl\_chat} from \texttt{icl\_flat} in §2.3.

What is missing is the comparison itself: four modern preference-optimisation losses (DPO, IPO, SimPO, KTO), two ICL variants, the same Llama-3-8B-Instruct base, the same per-user preference data, and a sweep of training sizes that brackets the cold-start regime end-to-end. That is the gap this thesis fills.